\section{Conclusion}
\label{sec:conclusion}

In this work, we revisited and analyzed two existing query optimization strategies.  With our experimental evaluation, we confirm the theoretical findings that \holistic induces a significant optimization overhead, especially for larger queries\revisionadd{, even when applying pruning}.  Nevertheless, we showcase scenarios where this overhead is justified for accelerating query execution and present an optimization guideline to detect such scenarios.  In three microbenchmarks, \holistic outperforms \split by up to one order of magnitude.  Furthermore, we propose a hybrid \topk optimization strategy and demonstrate that it is able to recover the holistic QEPs while improving the optimization time by multiple orders of magnitude compared to \holistic.  Additionally, we survey the TPC-H, JOB, and CEB benchmarks, where \split is superior for the majority of the queries.  However, we identify some queries for which \holistic and especially \topk are able to outperform \split by approximately \revisionrewrite{1.4x}.  \revisionadd{Lastly, we propose a simple guideline that robustly determines a suitable optimization approach to apply for a given query.}
